\documentclass[aspectratio=169, 10pt]{beamer}
\usepackage{sunass_estilo}

\title{Módulo X: Finanzas y Evaluación de Proyectos}
\subtitle{Guía Teórica de Repaso Rápido · PEU SUNASS}
\author[PEU SUNASS · Economía]{Preparación Examen de Admisión PEU SUNASS}
\institute{Basado en el Temario Oficial de Economía · Finanzas Corporativas y Regulatorias}
\date[2026]{2026}

\begin{document}

\begin{frame}[plain]
  \titlepage
\end{frame}

% ==============================================================================
% SLIDE 2: MAPA CONCEPTUAL Y VALOR DEL DINERO EN EL TIEMPO
% ==============================================================================
\begin{frame}{1. Capital, Tasas de Interés y Valor del Dinero en el Tiempo}
  \vspace{0.1cm}
  \begin{columns}[T]
    \begin{column}{0.32\textwidth}
      \tarjeta{azulPrincipal}{Capitalización y TEA}{
        • \textbf{Compuesta:} \(VF = VP(1+i)^n\).\\
        • \textbf{Tasa Efectiva (TEA):} \(TEA = \left(1 + \frac{TNA}{m}\right)^m - 1\).\\
        • \textbf{Continua:} \(VF = VP \cdot e^{rn}\).
      }
    \end{column}
    \begin{column}{0.32\textwidth}
      \tarjeta{naranjaAcento}{Rentas y Perpetuidades}{
        • \textbf{Perpetuidad:} \(VP = \frac{C}{r}\) (con crecimiento \(g < r\): \(VP = \frac{C_1}{r - g}\)).\\
        • \textbf{Anualidad vencida:} \(VP = C \cdot \frac{1 - (1+r)^{-n}}{r}\).
      }
    \end{column}
    \begin{column}{0.32\textwidth}
      \tarjeta{verdeAcento}{Ecuación de Fisher (Inflación)}{
        \((1 + i_{\text{nom}}) = (1 + r_{\text{real}})(1 + \pi)\)\\[3pt]
        • Aproximación: \(r \approx i - \pi\).\\
        • \textbf{Consistencia:} flujos nominales se descuentan con tasa nominal; flujos reales con tasa real.
      }
    \end{column}
  \end{columns}

  \vspace{0.15cm}
  \recuadroclave{
    \small\textbf{Regla de Oro de Consistencia:} Nunca mezcles flujos de caja en soles constantes (reales) con una tasa de descuento nominal (que incluye inflación).
  }
  \mensajefinal{Un sol hoy vale más que un sol mañana por su costo de oportunidad (tasa de interés), riesgo e inflación.}
\end{frame}

% ==============================================================================
% SLIDE 3: INDICADORES FINANCIEROS Y DUPONT
% ==============================================================================
\begin{frame}{2. Interpretación de Principales Indicadores Financieros}
  \vspace{0.1cm}
  \begin{columns}[T]
    \begin{column}{0.32\textwidth}
      \tarjeta{azulPrincipal}{Liquidez (Corto Plazo)}{
        • \textbf{Liquidez Corriente:} \(\frac{\text{Activo Cte.}}{\text{Pasivo Cte.}}\).\\
        • \textbf{Prueba Ácida:} \(\frac{\text{Activo Cte.} - \text{Inventarios}}{\text{Pasivo Cte.}}\).\\
        • \textbf{Capital de Trabajo:} \(AC - PC\).
      }
    \end{column}
    \begin{column}{0.32\textwidth}
      \tarjeta{naranjaAcento}{Solvencia y Cobertura}{
        • \textbf{Endeudamiento:} \(\frac{\text{Pasivo Total}}{\text{Patrimonio}}\) o \(\frac{D}{A}\).\\
        • \textbf{Cobertura de Intereses:} \(\frac{EBIT}{\text{Gastos Financieros}}\).\\
        • \textbf{EBITDA:} Utilidad operativa + Depreciación + Amortización (proxy de flujo operativo).
      }
    \end{column}
    \begin{column}{0.32\textwidth}
      \tarjeta{verdeAcento}{Rentabilidad e Identidad DuPont}{
        • \(ROA = \frac{\text{Utilidad Neta}}{\text{Activos}}\).\\
        • \textbf{Identidad DuPont:}\\
        \(ROE = \underbrace{\frac{UN}{\text{Ventas}}}_{\text{Margen}} \times \underbrace{\frac{\text{Ventas}}{\text{Activos}}}_{\text{Rotación}} \times \underbrace{\frac{\text{Activos}}{\text{Patrimonio}}}_{\text{Apalancamiento}}\)
      }
    \end{column}
  \end{columns}

  \vspace{0.15cm}
  \recuadroclave{
    \small\textbf{Flujo de Caja Libre (FCL) vs. Accionista (FCA):} \quad \(FCL = EBIT(1-T) + \text{Dep} - CAPEX - \Delta CTN\) (se descuenta al \textbf{WACC}; el \(FCA\) se descuenta al \(k_e\)).
  }
  \mensajefinal{Un proyecto puede mostrar utilidad contable positiva y aun así tener flujo de caja negativo por inversiones en CAPEX o capital de trabajo.}
\end{frame}

% ==============================================================================
% SLIDE 4: CRITERIOS DE EVALUACIÓN (VAN, TIR, CAE)
% ==============================================================================
\begin{frame}{3. Evaluación de Proyectos: VAN, TIR, B/C y CAE}
  \begin{columns}[T]
    \begin{column}{0.52\textwidth}
      {\small\textbf{Perfil del VAN y Conflicto con la TIR}}
      \vspace{0.05cm}
      \begin{tikzpicture}[scale=0.78]
        \begin{axis}[
          axis lines=left,
          xlabel={Tasa de descuento ($r$)}, ylabel={$VAN(r)$},
          xmin=0, xmax=0.35, ymin=-30, ymax=110,
          xtick={0.10, 0.16, 0.25}, xticklabels={$r^*$, $r_{\text{Fisher}}$, $TIR_A$},
          ytick={0},
          width=8.0cm, height=5.1cm
        ]
          % Eje 0
          \draw[grisTexto, thin] (axis cs:0,0) -- (axis cs:0.35,0);
          % Proyecto A (flujos tardios -> mayor pendiente)
          \addplot[azulPrincipal, very thick, domain=0:0.30] {100 - 400*x};
          \node[azulPrincipal, font=\scriptsize, anchor=west] at (axis cs:0.02, 92) {Proy. A (Mayor escala)};
          % Proyecto B (flujos tempranos -> mayor TIR, menor VAN a r*)
          \addplot[rojoAcento, very thick, domain=0:0.34] {68 - 200*x};
          \node[rojoAcento, font=\scriptsize, anchor=west] at (axis cs:0.18, 36) {Proy. B (Mayor TIR)};
          % Tasa de Fisher (100-400x = 68-200x => 32 = 200x => x = 0.16)
          \draw[dotted, verdeAcento, thick] (axis cs:0.16, 0) -- (axis cs:0.16, 36);
          \draw[dashed, grisTexto] (axis cs:0.10, 0) -- (axis cs:0.10, 60);
          \addplot[only marks, mark=*, mark size=2.4pt, verdeAcento] coordinates {(0.16, 36)};
        \end{axis}
      \end{tikzpicture}
    \end{column}

    \begin{column}{0.45\textwidth}
      \indicador{Valor Actual Neto (VAN)}{azulPrincipal}{\(VAN = -I_0 + \sum_{t=1}^n \frac{FC_t}{(1+r)^t} > 0\)}
      {\scriptsize
        • Mide la creación absoluta de riqueza. \textbf{Criterio supremo} ante proyectos mutuamente excluyentes.\\[4pt]
      }
      \indicador{Limitaciones de la TIR}{rojoAcento}{\text{Falla si hay escala distinta o flujos no convencionales}}
      {\scriptsize
        • \textbf{Múltiples TIR:} si los flujos cambian de signo más de una vez (Regla de Descartes).\\
        • \textbf{Tasa de Fisher:} tasa donde \(VAN_A = VAN_B\). Si \(r < r_{\text{Fisher}}\), la TIR puede contradecir al VAN.\\
        • \textbf{Vidas útiles distintas:} usar \textbf{CAE} (Costo Anual Equivalente) o repetir cadena.
      }
    \end{column}
  \end{columns}

  \mensajefinal{Ante cualquier conflicto entre TIR y VAN en proyectos mutuamente excluyentes, siempre se elige el proyecto de mayor VAN.}
\end{frame}

% ==============================================================================
% SLIDE 5: RIESGO, CAPM Y WACC REGULATORIO (SUNASS)
% ==============================================================================
\begin{frame}{4. Riesgo, CAPM y WACC Regulatorio en EPS (SUNASS)}
  \vspace{0.1cm}
  \begin{columns}[T]
    \begin{column}{0.32\textwidth}
      \tarjeta{azulPrincipal}{CAPM y Costo del Patrimonio}{
        \(k_e = r_f + \beta_L (E[R_m] - r_f) + RP\)\\[4pt]
        • Solo se remunera el \textbf{riesgo sistemático} (\(\beta\)), no el riesgo diversificable.\\
        • \(RP\): Prima por riesgo país (EMBI Peru).
      }
    \end{column}
    \begin{column}{0.32\textwidth}
      \tarjeta{naranjaAcento}{Apalancamiento (Hamada)}{
        \(\beta_L = \beta_U \left[1 + (1 - T)\frac{D}{E}\right]\)\\[4pt]
        • \(\beta_U\): beta del negocio puro (desapalancado).\\
        • A mayor deuda \(D/E\), sube el riesgo financiero y sube \(k_e\).
      }
    \end{column}
    \begin{column}{0.32\textwidth}
      \tarjeta{verdeAcento}{WACC Regulatorio SUNASS}{
        \(WACC = \frac{E}{V}k_e + \frac{D}{V}k_d(1 - T)\)\\[4pt]
        • Tasa de descuento oficial para calcular la \textbf{Base de Capital} y la tarifa media en el Plan Maestro Optimizado (PMO).\\
        • Incluye el \textbf{escudo fiscal} \((1-T)\) de la deuda.
      }
    \end{column}
  \end{columns}

  \vspace{0.15cm}
  \recuadroclave{
    \small\textbf{Fórmula de Tarifa Media (SUNASS):} El Valor Presente de los Ingresos Tarifarios debe cubrir el VP de Costos de O\&M + Inversiones (CAPEX) + Impuestos, descontados al \textbf{WACC real} aprobado por el regulador.
  }
  \mensajefinal{En el modelo CAPM, el riesgo no sistemático se elimina diversificando el portafolio y no recibe prima de retorno.}
\end{frame}

% ==============================================================================
% SLIDE 6: VALUACIÓN DE BONOS, DURACIÓN Y ACCIONES (GORDON-SHAPIRO)
% ==============================================================================
\begin{frame}{5. Valuación de Bonos, Duración y Acciones}
  \vspace{0.05cm}
  \begin{columns}[T]
    \begin{column}{0.32\textwidth}
      \tarjeta{azulPrincipal}{Precio de Bonos y YTM}{
        \(P_B = \sum_{t=1}^n \frac{C}{(1+y)^t} + \frac{VN}{(1+y)^n}\)\\[3pt]
        • \textbf{Relación inversa} entre precio \(P_B\) y tasa de rendimiento (\(YTM = y\)).\\
        • Si \(c = y \Rightarrow\) Bono \textbf{a la par} (\(P=VN\)).\\
        • Si \(c > y \Rightarrow\) \textbf{con prima}; si \(c < y \Rightarrow\) \textbf{bajo par}.
      }
    \end{column}
    \begin{column}{0.32\textwidth}
      \tarjeta{naranjaAcento}{Duración y Convexidad}{
        • \textbf{Duración de Macaulay (\(D\)):} vida media ponderada de los flujos.\\
        • \textbf{Duración Modificada:} \(D_{\text{mod}} = \frac{D}{1+y}\).\\
        • Aproximación lineal: \(\frac{\Delta P}{P} \approx -D_{\text{mod}}\,\Delta y\) (a mayor vencimiento y menor cupón, mayor riesgo de tasa).
      }
    \end{column}
    \begin{column}{0.32\textwidth}
      \tarjeta{verdeAcento}{Acciones (Gordon-Shapiro)}{
        Dividendos creciendo a tasa constante \(g < k_e\):\\[3pt]
        \(P_0 = \frac{D_1}{k_e - g} = \frac{D_0(1+g)}{k_e - g}\)\\[3pt]
        • \(g = ROE \times b\) (tasa de retención de utilidades).\\
        • Si el proyecto invierte al \(ROE = k_e\), el PVGO es cero.
      }
    \end{column}
  \end{columns}

  \vspace{0.08cm}
  \recuadroclave{
    \small\textbf{Racionamiento de Capital y TIRM:} Con presupuesto limitado se ordenan proyectos por el \textbf{Índice de Rentabilidad} (\(IR = \frac{VP(\text{Flujos})}{I_0} > 1\)); la \textbf{TIRM} corrige a la TIR asumiendo reinversión al \(WACC\) y no a la propia TIR.
  }
  \mensajefinal{Los bonos cupón cero tienen una Duración de Macaulay exactamente igual a su plazo de vencimiento (\(D = n\)).}
\end{frame}

% ==============================================================================
% SLIDE 7: MARKOWITZ (CML VS SML), MODIGLIANI-MILLER Y OPCIONES
% ==============================================================================
\begin{frame}{6. Portafolio (CML vs. SML), Modigliani–Miller y Opciones}
  \vspace{0.05cm}
  \begin{columns}[T]
    \begin{column}{0.32\textwidth}
      \tarjeta{azulPrincipal}{Markowitz: CML vs. SML}{
        • \textbf{Ratio de Sharpe:} \(SR = \frac{E[R_p] - r_f}{\sigma_p}\).\\
        • \textbf{CML (Línea del Mercado de Capitales):} relaciona \(E[R_p]\) con riesgo total \(\sigma_p\) (solo para \textbf{portafolios eficientes}).\\
        • \textbf{SML (Línea del Mercado de Valores):} relaciona \(E[R_i]\) con \(\beta_i\) (para \textbf{todo activo individual}).
      }
    \end{column}
    \begin{column}{0.32\textwidth}
      \tarjeta{naranjaAcento}{Modigliani–Miller (MM I y II)}{
        • \textbf{Sin impuestos:} la estructura de capital es \textbf{irrelevante} (\(V_L = V_U\)); el \(WACC\) es constante y \(k_e\) sube linealmente con \(D/E\).\\
        • \textbf{Con impuestos:} \(V_L = V_U + T\cdot D\) (valor del escudo fiscal).\\
        • \textbf{Trade-off:} equilibra escudo fiscal vs. costos de quiebra.
      }
    \end{column}
    \begin{column}{0.32\textwidth}
      \tarjeta{verdeAcento}{Derivados y Paridad Put–Call}{
        • \textbf{Call} (derecho a comprar) vs. \textbf{Put} (derecho a vender) al precio de ejercicio \(K\).\\
        • \textbf{Paridad Put–Call (europeas):}\\
        \(\quad C + \frac{K}{(1+r)^T} = P + S_0\)\\
        • \textbf{Forwards / Futuros:} fijan tipo de cambio o precio futuro para cubrir riesgos.
      }
    \end{column}
  \end{columns}

  \vspace{0.08cm}
  \recuadroclave{
    \small\textbf{Alpha de Jensen (\(\alpha_i\)):} Si un activo está por \textbf{encima} de la SML (\(\alpha_i = R_i - [r_f + \beta_i(R_m - r_f)] > 0\)), está \textbf{subvaluado} y conviene comprarlo.
  }
  \mensajefinal{En mercados eficientes sin fricciones, las decisiones de financiamiento no crean valor; solo lo hacen las decisiones de inversión con \(VAN > 0\).}
\end{frame}

% ==============================================================================
% SLIDE 8: CONCLUSIONES Y CHEAT-SHEET MÓDULO X
% ==============================================================================
\begin{frame}{Conclusiones y Formulario Rápido — Módulo X}
  \vspace{0.05cm}
  \recuadroclave{
    \small\textbf{Regla de Oro:} Todo activo vale el valor presente de sus flujos futuros descontados a una tasa coherente con su riesgo (WACC / CAPM) y su unidad monetaria (nominal vs. real).
  }

  \vspace{0.15cm}
  \begin{columns}[T]
    \begin{column}{0.32\textwidth}
      \tarjeta{verdeAcento}{Fórmulas Imprescindibles}{
        • Fisher: \((1+i) = (1+r)(1+\pi)\)\\
        • CAPM: \(k_e = r_f + \beta_L(R_m - r_f)\)\\
        • WACC: \(\frac{E}{V}k_e + \frac{D}{V}k_d(1-T)\)
      }
    \end{column}
    \begin{column}{0.32\textwidth}
      \tarjeta{azulPrincipal}{Trampas de Examen}{
        • Usar TIR en vez de VAN cuando hay distinta escala o flujos no convencionales.\\
        • Comparar VAN directo en proyectos con diferente vida útil (usar CAE).\\
        • Olvidar el escudo fiscal \((1-T)\) en \(k_d\).
      }
    \end{column}
    \begin{column}{0.32\textwidth}
      \tarjeta{rojoAcento}{Siguientes Pasos}{
        • Practicar Bloques B y C del Módulo 10 en el Simulador.\\
        • Realizar un Simulacro Completo de 90 preguntas (150 min).
      }
    \end{column}
  \end{columns}

  \vspace{0.25cm}
  \begin{center}
    {\color{azulPrincipal}\Large\textbf{¡Éxito!}} \quad {\color{grisTexto}Módulo X · Finanzas y Evaluación de Proyectos · PEU SUNASS 2026}
  \end{center}
\end{frame}

\end{document}
